\chapter{Diskussion und Ausblick}\label{sec:fazit}


% - Es konnte folgendes allgemein herausgefunden werden:
%  -> Starke Unterdrückung bei halber Füllung -> Quantum kritischer Punkt, was auch mit anderen Arbeiten konsistent ist.
%  -> s-wave SC am besten bei VHS (1t)

% - Es konnte folgendes über Graphen heruasgefunden werden:
%  -> Keine s-wave Supraleitung bei mu = 0 -> Eperimentell auch nicht gefunden (Zitierung schwierig?)
%  -> Bei dotiertem Graphen ist Supraleitung möglich, Vorhersage auch experimentell bestätigt (Li-decorated SLG) Delta = 0.9\,meV
%  -> U kann aber als Parameter nur eingeschränkt werden, genaue Vorhersagen schwierig.
%  -> s-wave könnte da also einen relevanten Beitrag zu bringen, obwohl Anistrope Gap (extended s-wave?) bzw. proposed d-wave bei heavy doping (ZITAT Topological Phases)

% Ausblick:
% - Nicht beachtet wurde die Veränderung der Bänder bei Heavy-Doping (ZITAT Paper)
% - Experimentelle Werte könnten U festlegen -> genauere Vorhersage 
% - Modell bedingt auf Transition Metals wie MoS2 anwendbar (mit $W = \ldots eV, E_D = \ldots meV$), aber 1-Band TB zu einfach (ZITAT 11Band TB)

In dieser Arbeit wurde 
